\clearpage
\section{Data Formats}

\subsection{Integer Data Sizes}
Integer operands use size suffixes to select the low-order subfield of an Rn register and the number of bytes
transferred by memory operands.

Byte, word, and long operations use the low-order subfield of the selected register. Q-sized operations use the full
64-bit register.

An Rn destination write smaller than Q replaces only the selected low-order subfield; bits above the operation size
are preserved. Q-sized Rn writes replace all 64 bits. Memory destination writes transfer only the selected number of
bytes.

An instruction may explicitly define a full-register write, zero-extension, sign-extension, or another upper-bit
rule. Such instruction-specific rules override the default subfield write rule.

\manualtablecaption{Scalar Data Sizes}
\begin{manualtabular}{@{}>{\raggedright\arraybackslash}p{0.45in}>{\raggedright\arraybackslash}p{0.55in}>{\raggedright\arraybackslash}p{0.45in}>{\raggedright\arraybackslash}p{0.50in}>{\raggedright\arraybackslash}p{1.55in}@{}}
\toprule
\rowcolor{ManualHeaderFill}
\textbf{Code} & \textbf{Suffix} & \textbf{Bits} & \textbf{Bytes} & \textbf{Name}\\
\midrule
\texttt{B} & \texttt{@B_SUFFIX@} & @B_BITS@ & @B_BYTES@ & Byte\\
\texttt{W} & \texttt{@W_SUFFIX@} & @W_BITS@ & @W_BYTES@ & Word\\
\texttt{L} & \texttt{@L_SUFFIX@} & @L_BITS@ & @L_BYTES@ & Long\\
\texttt{Q} & \texttt{@Q_SUFFIX@} & @Q_BITS@ & @Q_BYTES@ & Quad\\
\texttt{S} & \texttt{@S_SUFFIX@} & @S_BITS@ & @S_BYTES@ & Single\\
\texttt{D} & \texttt{@D_SUFFIX@} & @D_BITS@ & @D_BYTES@ & Double\\
\bottomrule
\end{manualtabular}

\begin{manuallistedformatdiagram}{Integer Register Operand Subfields}{4}
\manualformatrow{Q}{%
\manualformatfield{Q operand bits 63..0}{64}
}
\manualformatrow{L}{%
\manualformatfield{unchanged / instruction-defined}{32}
\manualformatfield{L operand bits 31..0}{32}
}
\manualformatrow{W}{%
\manualformatfield{unchanged / instruction-defined}{48}
\manualformatfield{W bits 15..0}{16}
}
\manualformatrow{B}{%
\manualformatfield{unchanged / instruction-defined}{56}
\manualformatfield{B bits 7..0}{8}
}
\end{manuallistedformatdiagram}


\subsection{Little-Endian Memory Organization}
Integer data, immediates, displacements, and absolute address payloads are little-endian. Multi-byte numeric
payloads are stored least significant byte first. The instruction header is byte-coded in instruction-stream
order rather than interpreted as a little-endian integer.

For an N-byte integer value stored at byte address A, byte A contains bits 7..0, byte A+1 contains bits 15..8, and
so on.

\begin{manuallistedbyteorderdiagram}{Little-Endian Byte Order for a 64-Bit Value}{byte}{address}{increasing addresses}{least significant byte first}{most significant byte last}
\manualbytecell{$V[7..0]$}{A+0}
\manualbytecell{$V[15..8]$}{A+1}
\manualbytecell{$V[23..16]$}{A+2}
\manualbytecell{$V[31..24]$}{A+3}
\manualbytecell{$V[39..32]$}{A+4}
\manualbytecell{$V[47..40]$}{A+5}
\manualbytecell{$V[55..48]$}{A+6}
\manualbytecell{$V[63..56]$}{A+7}
\end{manuallistedbyteorderdiagram}

\begin{manuallistedbyteorderdiagram}{Instruction Start Bytes}{header byte}{}{instruction stream byte order}{}{}
\manualbytecell{\shortstack{class/length\\opcode high}}{PC+0}
\manualbytecell{\shortstack{opcode bits\\if present}}{PC+1}
\end{manuallistedbyteorderdiagram}


\subsection{Instruction Payload Ordering}
The instruction header occupies one byte for extrashort forms, two bytes for short forms, and at least two start
bytes for extended forms. If a descriptor, immediate, displacement, or bitmap follows, those payload bytes
occupy increasing addresses in decode order.

Signed displacement and immediate payloads use two's-complement representation at their declared width before any
sign extension performed by the consuming instruction or EA form.

\manualtablecaption{Immediate, Displacement, and Address Payload Order}
\begin{manuallongtable}{@{}>{\raggedright\arraybackslash}p{0.75in}>{\raggedright\arraybackslash}p{0.45in}>{\raggedright\arraybackslash}p{4.20in}@{}}
\toprule
\rowcolor{ManualHeaderFill}
\textbf{Payload} & \textbf{Bytes} & \textbf{Memory Order}\\
\midrule
\endfirsthead
\multicolumn{3}{l}{\scriptsize\itshape Table \themanualtable\ (continued)}\\
\toprule
\rowcolor{ManualHeaderFill}
\textbf{Payload} & \textbf{Bytes} & \textbf{Memory Order}\\
\midrule
\endhead
\texttt{imm8} & 1 & single byte\\
\texttt{imm16} & 2 & bits 7..0 at the lower byte, then bits 15..8\\
\texttt{imm32} & 4 & least significant byte first\\
\texttt{imm64} & 8 & least significant byte first\\
\texttt{disp8} & 1 & single signed two's-complement byte\\
\texttt{disp16} & 2 & signed two's-complement value, little-endian\\
\texttt{disp32} & 4 & signed value, least significant byte first\\
\texttt{disp64} & 8 & declared-width payload, least significant byte first\\
\texttt{abs32} & 4 & absolute-address payload, least significant byte first\\
\texttt{abs64} & 8 & absolute-address payload, least significant byte first\\
\bottomrule
\end{manuallongtable}

\manualtablecaption{Immediate Operand Interpretation}
\begin{manuallongtable}{@{}>{\raggedright\arraybackslash}p{0.65in}>{\raggedright\arraybackslash}p{0.38in}>{\raggedright\arraybackslash}p{0.62in}>{\raggedright\arraybackslash}p{0.62in}>{\raggedright\arraybackslash}p{0.90in}>{\raggedright\arraybackslash}p{2.00in}@{}}
\toprule
\rowcolor{ManualHeaderFill}
\textbf{Operand} & \textbf{Bits} & \textbf{Value} & \textbf{Range} & \textbf{Extension} & \textbf{Applies When}\\
\midrule
\endfirsthead
\multicolumn{6}{l}{\scriptsize\itshape Table \themanualtable\ (continued)}\\
\toprule
\rowcolor{ManualHeaderFill}
\textbf{Operand} & \textbf{Bits} & \textbf{Value} & \textbf{Range} & \textbf{Extension} & \textbf{Applies When}\\
\midrule
\endhead
\texttt{imm3} & @IMM3_WIDTH@ & unsigned & @IMM3_RANGE@ & zero-extend & PUSHP/POPP canonical register-pair selector\\
\texttt{imm4} & @IMM4_WIDTH@ & unsigned & @IMM4_RANGE@ & zero-extend & SETF/CLRF integer FLAGS masks\\
\texttt{imm6} & @IMM6_WIDTH@ & unsigned & @IMM6_RANGE@ & zero-extend & integer forms with an explicit B/W/L/Q operation size\\
\texttt{imm7} & @IMM7_WIDTH@ & unsigned & @IMM7_RANGE@ & zero-extend & EXTRACT register-pair forms\\
\bottomrule
\end{manuallongtable}

PUSHP and POPP interpret \texttt{imm3} as a canonical register-pair index, so every encoding from 0 through 7 is valid.
SETF and CLRF interpret \texttt{imm4} as the low FLAGS nibble: bit 3 selects Z, bit 2 selects N, bit 1 selects C,
and bit 0 selects V.

An \texttt{imm6} value is zero-extended to the selected operation size before use when that size is wider than six
bits. EXTRACT uses \texttt{imm7} as the offset into its concatenated register pair.

\subsection{Floating-Point Data Formats}
Floating-point scalar formats are little-endian in memory as well. The S format is 32 bits and the D format is 64
bits.

The bit-level floating-point encoding, NaN policy, exception flags, and rounding behavior are defined by the
floating-point instruction and state descriptions; their memory byte order follows the same
least-significant-byte-first rule as integer data.
